\documentclass[10pt,a4paper]{article}
\usepackage[margin=18mm,headheight=14pt]{geometry}
\usepackage{fontspec}
\setmainfont{lmroman10-regular.otf}[BoldFont=lmroman10-bold.otf,ItalicFont=lmroman10-italic.otf,BoldItalicFont=lmroman10-bolditalic.otf]
\setsansfont{lmsans10-regular.otf}[BoldFont=lmsans10-bold.otf]
\setmonofont{lmmonolt10-regular.otf}[BoldFont=lmmonolt10-bold.otf,ItalicFont=lmmonolt10-oblique.otf]
\usepackage[ngerman]{babel}
\usepackage{amsmath,amssymb,booktabs,tabularx,array,fancyhdr,enumitem}
\usepackage[hidelinks]{hyperref}
\usepackage{tikz}
\usetikzlibrary{shapes.gates.logic.US,positioning,calc}
\pagestyle{fancy}\fancyhf{}
\fancyhead[L]{\sffamily AI-Hardware · Tag 5}
\fancyhead[R]{\sffamily Abschluss · Lösungen}
\fancyfoot[L]{\small Erst selbst bearbeiten, dann vergleichen}
\fancyfoot[R]{\thepage\ / 4}
\setlength{\emergencystretch}{3em}\setlength{\parindent}{0pt}\setlength{\parskip}{5pt}
\setlength{\tabcolsep}{7pt}\renewcommand{\arraystretch}{1.18}
\newcommand{\titlepagepart}[2]{\vspace*{-4mm}{\sffamily\small #1}\par{\sffamily\LARGE\bfseries #2}\par\vspace{2mm}}
\newcommand{\block}[1]{\par\vspace{2mm}{\sffamily\large\bfseries #1}\par}
\newcommand{\code}[1]{\texttt{#1}}
\newcommand{\N}{\mathrm{N}}
\begin{document}

\titlepagepart{01 / LÖSUNGEN ZU SEITE 1 UND 2}{Boolesche Algebra und De Morgan}
Wo mehrere Wege zum Ziel führen, ist einer ausgeführt; ein anderer, korrekt
begründeter Weg ist ebenso gültig. Entscheidend ist, dass jeder Schritt ein
benanntes Gesetz benutzt.

\block{1. Vereinfachen mit den Gesetzen}
\textbf{a)}
\begin{align*}
Y&=\overline{A}B+AB+A\overline{B}\\
 &=B(\overline{A}+A)+A\overline{B} &&\text{Distributivgesetz}\\
 &=B\cdot 1+A\overline{B} &&\text{Komplement}\\
 &=B+A\overline{B} &&\text{Neutrales Element}\\
 &=A+B &&\text{Form } X+\overline{X}Y=X+Y \text{ mit } X=B,\ Y=A
\end{align*}
Die letzte Zeile ausführlich: $B+A\overline{B}=(B+A)(B+\overline{B})=(B+A)\cdot 1=A+B$.

\textbf{b)}
$(A+B)(A+\overline{B})\overset{\text{Distributiv}}{=}A+B\overline{B}
\overset{\text{Komplement}}{=}A+0\overset{\text{Neutral}}{=}A.$
Hier wird die OR-Form $A+BC=(A+B)(A+C)$ von rechts nach links gelesen. Die
häufigste Falle ist, den Ausdruck wie gewöhnliches Ausmultiplizieren zu behandeln.

\textbf{c)}
$\overline{\overline{A}+\overline{B}}
\overset{\text{De Morgan}}{=}\overline{\overline{A}}\cdot\overline{\overline{B}}
\overset{\text{doppelte Negation}}{=}A\,B.$

\textbf{d)} Es gilt $AB+\overline{A}C+BC=AB+\overline{A}C$; der Term $BC$ entfällt.
$BC$ ist nur dann 1, wenn $B=C=1$. Ist zusätzlich $A=1$, liefert bereits $AB$
eine 1; ist $A=0$, liefert bereits $\overline{A}C$ eine 1. $BC$ erzeugt also in
keiner Zeile eine 1, die nicht ohnehin entsteht.

\block{2. De Morgan sicher anwenden}
\textbf{a)}
\[
\overline{A(B+\overline{C})}
=\overline{A}+\overline{(B+\overline{C})}
=\overline{A}+\overline{B}\cdot\overline{\overline{C}}
=\overline{A}+\overline{B}\,C .
\]
Der Überstrich wird über \emph{jede} Verknüpfung einzeln gezogen, außen beginnend.

\textbf{b)} Falsch ist die Annahme, De Morgan negiere lediglich alle Variablen.
Tatsächlich wechselt auf \emph{jeder} Ebene auch die Verknüpfung:
\[
\overline{A+BC}=\overline{A}\cdot\overline{BC}
=\overline{A}\,(\overline{B}+\overline{C})
=\overline{A}\,\overline{B}+\overline{A}\,\overline{C}.
\]
Gegenbeispiel $A=0$, $B=1$, $C=1$: links $\overline{0+1}=0$, in der falschen
Form $1+0+0=1$.

\textbf{c)} Wird die Ausgangsblase über den Gatterkörper geschoben, tauscht der
Gattertyp zwischen AND und OR, und an \emph{jedem} Eingang entsteht eine Blase.
Aus NAND wird damit ein OR mit zwei Eingangsblasen.

\block{3. Die NAND-Aufbauten als Gleichung}
\begin{center}
\begin{tabular}{llll}\toprule
Funktion&Zwischenwerte&Ergebnis&NANDs\\\midrule
NOT&---&$Y=\N(A,A)=\overline{AA}=\overline{A}$&1\\
AND&$T=\N(A,B)$&$Y=\N(T,T)=\overline{\overline{AB}}=AB$&2\\
OR&$P=\N(A,A),\ Q=\N(B,B)$&$Y=\N(P,Q)=\overline{\overline{A}\,\overline{B}}=A+B$&3\\\bottomrule
\end{tabular}
\end{center}
NOT trägt die Idempotenz $AA=A$, AND die doppelte Negation, OR De Morgan.

\newpage
\titlepagepart{02 / LÖSUNGEN ZU SEITE 3}{Netze rechnen}

\block{4. Ein unbekanntes Netz bestimmen}
\textbf{a)} Mit $P=\N(A,A)=\overline{A}$ und $Q=\N(A,B)=\overline{AB}$:
\begin{center}
\begin{tabular}{cc|ccc}\toprule
$A$&$B$&$P$&$Q$&$Y$\\\midrule
0&0&1&1&0\\
0&1&1&1&0\\
1&0&0&1&1\\
1&1&0&0&1\\\bottomrule
\end{tabular}
\end{center}
Die Ergebnisspalte ist identisch mit $A$; $B$ hat keinen Einfluss auf $Y$.

\textbf{b)}
\[
Y=\overline{P\,Q}=\overline{\overline{A}\cdot\overline{AB}}
\;\overset{\text{De Morgan}}{=}\;\overline{\overline{A}}+\overline{\overline{AB}}
\;\overset{\text{dopp. Negation}}{=}\;A+AB
\;\overset{\text{Absorption}}{=}\;A .
\]
Die tragenden Gesetze sind De Morgan (mit doppelter Negation) und die
Absorption $A+AB=A$.

\textbf{c)} Drei Gatter können logisch dasselbe leisten wie eine durchgehende
Leitung. Eine korrekte Wahrheitstabelle beweist nur die Funktion, nicht die
Sparsamkeit des Aufbaus. Vor dem Bauen lohnt daher die algebraische
Vereinfachung: überflüssige Gatter kosten Fläche und Energie, ohne etwas
beizutragen. Umgekehrt gilt: Die Gatterzahl ist kein Maß für die Funktion.

\block{5. Halbaddierer: die Brücke zu Tag 4}
\textbf{a)}
\begin{center}
\begin{tabular}{cc|cc}\toprule
$A$&$B$&$S$&$C_{\text{out}}$\\\midrule
0&0&0&0\\
0&1&1&0\\
1&0&1&0\\
1&1&0&1\\\bottomrule
\end{tabular}
\end{center}
In der letzten Zeile ist $C_{\text{out}}S=10_2=2$, also das korrekte Ergebnis
von $1+1$ in zwei Bits.

\textbf{b)} Im XOR-Netz ist $T=\N(A,B)=\overline{AB}$ bereits vorhanden.
Der Übertrag ist dessen Negation:
\[
C_{\text{out}}=\N(T,T)=\overline{\overline{AB}}=AB .
\]
Mit Mitbenutzung von $T$: $4+1=5$ NAND-Gatter für den ganzen Halbaddierer.
Ohne Mitbenutzung, also XOR und AND getrennt aufgebaut: $4+2=6$ NAND-Gatter.

\textbf{c)} An der Wahrheitstabelle ändert die Mitbenutzung nichts; die
Funktion ist identisch. In der realen Schaltung treibt der Ausgang des ersten
NAND nun drei statt zwei Eingänge. Der Fan-out und damit die kapazitive Last
steigen, das Signal wird an dieser Stelle langsamer. Genau diesen Unterschied
zeigt eine reine Logiksimulation nicht.

\newpage
\titlepagepart{03 / LÖSUNGEN ZU SEITE 4}{Pegel, Störabstand, Kennlinie}

\block{6. Störabstände berechnen}
\textbf{a)}
\[
NM_L=V_{IL}-V_{OL}=0{,}99-0{,}33=0{,}66\ \text{V},\qquad
NM_H=V_{OH}-V_{IH}=2{,}97-1{,}98=0{,}99\ \text{V}.
\]
Beide Störabstände sind positiv, die Verbindung ist zulässig. Im ungünstigsten
Fall darf ein LOW um bis zu 0,66~V angehoben und ein HIGH um bis zu 0,99~V
abgesenkt werden, bevor der garantierte Bereich verlassen wird.

\textbf{b)}
\[
NM_L=0{,}40-0{,}50=-0{,}10\ \text{V},\qquad
NM_H=2{,}40-2{,}00=0{,}40\ \text{V}.
\]
$NM_L$ ist negativ: Der Treiber darf im schlechtesten Fall bis 0,50~V ausgeben,
der Empfänger akzeptiert aber nur bis 0,40~V sicher als LOW. Ein korrekt
gesendetes LOW kann damit als unbestimmt oder sogar als HIGH gelesen werden,
und zwar ohne jede äußere Störung. Die HIGH-Seite bleibt mit 0,40~V Reserve in
Ordnung; die Paarung ist in dieser Richtung dennoch unbrauchbar.

\textbf{c)} Der Störabstand beschreibt nur die \emph{untere} Schranke für ein
erkanntes HIGH. Zusätzlich gibt es eine \emph{obere} Grenze der zulässigen
Eingangsspannung, die sich aus der Bauteilphysik ergibt (Schutzstrukturen,
Gate-Oxid). 5,0~V an einem 3,3-V-Eingang ist logisch ein eindeutiges HIGH,
überschreitet aber diese Grenze und kann den Eingang dauerhaft schädigen.
Logische Gültigkeit und elektrische Zulässigkeit sind zwei getrennte Prüfungen.
Nötig wäre ein Pegelwandler oder ein ausdrücklich 5-V-toleranter Eingang.

\block{7. Von der Kennlinie zur Regel}
\textbf{a)}
\begin{center}
\begin{tikzpicture}[x=1cm,y=1cm]
\draw[->](0,0)--(5.2,0) node[right] {$V_{in}$};
\draw[->](0,0)--(0,3.6) node[above] {$V_{out}$};
\draw[thick](0,3).. controls (1.1,3) and (1.5,2.9)..(1.9,2.3)
..controls(2.3,1.6)and(2.6,.5)..(3.1,.2)..controls(3.5,.05)and(4.2,0)..(4.8,0);
\node[left]at(0,3){$V_{DD}$};\node[below]at(4.8,0){$V_{DD}$};
\fill(1.9,2.3)circle(2pt);\fill(3.1,.2)circle(2pt);
\draw[dashed](1.9,2.3)--(1.9,0) node[below]{$V_{IL}$};
\draw[dashed](3.1,.2)--(3.1,0) node[below]{$V_{IH}$};
\draw[dashed](1.9,2.3)--(0,2.3) node[left]{$V_{OH}$};
\draw[dashed](3.1,.2)--(0,.2) node[left]{$V_{OL}$};
\node[right] at (2.2,1.5) {Steigung $-1$};
\end{tikzpicture}
\end{center}
An den beiden Punkten mit $\dfrac{dV_{out}}{dV_{in}}=-1$ werden die
Eingangsschwellen $V_{IL}$ und $V_{IH}$ festgelegt; die zugehörigen
Ausgangswerte liefern $V_{OH}$ und $V_{OL}$. Die Zeichnung ist schematisch und
nicht maßstäblich.

\textbf{b)} Static Discipline: Liegen an einem Gatter gültige, stabile
Eingangspegel an, so liefert es nach seiner Einschwingzeit gültige
Ausgangspegel. Sie macht ausdrücklich \emph{keine} Aussage über das Verhalten
\emph{während} eines Übergangs; Verzögerung, Flankenform und die dabei
durchlaufenen Zwischenspannungen sind nicht Teil dieser Regel. Ebenso wenig
sagt sie etwas über Ströme, Lasten oder Leistungsaufnahme.

\textbf{c)} Keine der drei Teilaufgaben lässt sich in Logisim-Evolution
nachstellen. Das Werkzeug kennt nur die logischen Werte 0 und 1 sowie Fehler-
und Undefiniert-Zustände, aber keine Spannungswerte, keine Eingangsschwellen
und keine Bauteilgrenzen. Für Aufgabe 6 und 7 ist ngspice das passende
Werkzeug; genau dafür steht der Samstagsblock mit dem CMOS-Inverter.

\newpage
\titlepagepart{04 / LÖSUNGEN ZU SEITE 5}{Praxis und Erfolgskontrolle}

\block{8. Die offenen NAND-Aufbauten}
\begin{center}
\begin{tabular}{cc|ccccc}\toprule
$A$&$B$&NOT ($Y=\overline{A}$)&AND&OR&XOR&NOR\\\midrule
0&0&1&0&0&0&1\\
0&1&\textemdash&0&1&1&0\\
1&0&0&0&1&1&0\\
1&1&\textemdash&1&1&0&0\\\bottomrule
\end{tabular}
\end{center}
Erwartete NAND-Anzahl: NOT 1, AND 2, OR 3, XOR 4, NOR 4. Für NOR wird der
NAND-basierte OR-Ausgang (3 Gatter) mit einem weiteren NAND invertiert.
Zusammen sind das die fünf Aufbauten der Erfolgskontrolle.

\textbf{Kontrollfrage:} Von links nach rechts vorgehen und zuerst den
Zwischenwert prüfen, dessen Eingänge bereits korrekt sind, in den Netzen dieses
Tages also $T$ beziehungsweise $P$ und $Q$. Ein Fehler pflanzt sich fort; der
erste falsche Knoten von links lokalisiert die Ursache, während der Ausgang nur
zeigt, \emph{dass} etwas falsch ist. Die häufigsten Ursachen sind eine Leitung,
die optisch anliegt, aber keinen Verbindungspunkt hat, und eine Pin-Breite
ungleich einem Bit.

\block{9. Die beiden De-Morgan-Nachweise}
\begin{center}
\begin{tabular}{cc|cc|cc}\toprule
&&\multicolumn{2}{c|}{$\overline{AB}$ gegen $\overline{A}+\overline{B}$}&\multicolumn{2}{c}{$\overline{A+B}$ gegen $\overline{A}\,\overline{B}$}\\
$A$&$B$&$Y_1$&$Y_2$&$Y_1$&$Y_2$\\\midrule
0&0&1&1&1&1\\
0&1&1&1&0&0\\
1&0&1&1&0&0\\
1&1&0&0&0&0\\\bottomrule
\end{tabular}
\end{center}
Zwei Boolesche Funktionen sind genau dann gleich, wenn sie für \emph{jede}
Belegung denselben Wert liefern. Bei zwei Eingängen gibt es $2^2=4$ Belegungen;
die vollständige Tabelle prüft alle Fälle und ist damit ein erschöpfender
Beweis. Eine einzelne übereinstimmende Zeile schließt nichts aus: Auch die
\emph{falsche} Gleichung $\overline{AB}=\overline{A}\,\overline{B}$ stimmt in
den Zeilen $A=B=0$ und $A=B=1$ und scheitert erst bei 01 und 10.

\block{10. Erfolgskontrolle}
Der Plan verlangt fünf NAND-Schaltungen mit korrekter Wahrheitstabelle als
\code{.circ} im Repository. Aus den vorhandenen Dateien sind XOR und NOR sowie
beide De-Morgan-Vergleiche abgedeckt; neu hinzu kommen NOT, AND und OR. Die
Sammeldatei \code{tag05\_nand.circ} führt alle fünf als getrennte
Teilschaltungen zusammen, damit ein einziges Artefakt die Kontrolle erfüllt.
Die Einzeldateien können daneben bestehen bleiben; das \code{README.md} sollte
dann festhalten, welche Datei die maßgebliche ist.

Der Ordner \code{01\_digital} ist derzeit noch nicht versioniert. Erst der
Commit macht die Erfolgskontrolle nachprüfbar; ohne ihn ist die Arbeit zwar
getan, aber nicht belegt.

\textbf{Abschlussnotiz:} Individuell. Die ehrliche Antwort auf die Frage, was
heute nur benutzt und nicht verstanden wurde, ist der eigentliche Wert des
Tages. Sie gehört in \code{LOG.md} und bestimmt, was am Samstag im
Konsolidierungsblock aufgearbeitet wird.

\end{document}
